% !PRINTED 33
\weberpage{33}
%<<CONT
Vielfaches von $a$. Das Element 1 hat den Grad 1. Alle cyklischen Gruppen sind commutativ.

Es gilt nun für jede Abel'sche Gruppe $S$ der Fundamentalsatz, dass sich immer ein System von Elementen $A, B, C \ldots$ von den Graden $a, b, c \ldots$ so auswählen lässt, dass sich in der Form
\[
\Theta = A^\alpha B^\beta C^\gamma \ldots
\]
jedes Element $\Theta$ von $S$, und jedes nur einmal, darstellen lässt, wenn $\alpha$ ein volles Restsystem nach dem Modul $a$, $\beta$ ein volles Restsystem nach dem Modul $b$, $\gamma$ ein volles Restsystem nach dem Modul $c$ etc. durchläuft.

Ein System von Elementen, das zu einer solchen Darstellung geeignet ist, heisst eine \gesperrt{Basis der Gruppe}, und wir können den zu beweisenden Satz demnach so aussprechen:

\gesperrt{I. Jede Abel'sche Gruppe lässt sich durch eine Basis darstellen.}

Der Beweis des Satzes gründet sich auf folgende Reihe von Schlüssen:

\gesperrt{1. Ist $n$ der Grad der Gruppe $S$, sind $A_1, A_2 \ldots A_n$ ihre Elemente und $a_1, a_2 \ldots a_n$ deren Grade; durchlaufen ferner $\alpha_1, \alpha_2 \ldots \alpha_n$ volle Restsysteme nach den Moduln $a_1, a_2 \ldots a_n$, so wird in der Form}
\weq{1}{\Theta = A_1^{\alpha_1} A_2^{\alpha_2} \ldots A_n^{\alpha_n}}
\gesperrt{jedes Element von $S$ und jedes gleich oft dargestellt.}

Dass man jedes Element in der Form (1) überhaupt erhält, sieht man unmittelbar; denn man erhält z.\,B. $A_1$, wenn man $\alpha_1 = 1, \alpha_2 = 0 \ldots \alpha_n = 0$ setzt.

Es ist noch zu beweisen, dass man jedes Element gleich oft erhält.

Wenn aber
\weq{2}{1 = A_1^{h_1} A_2^{h_2} \ldots A_n^{h_n}}
eine Darstellung des Elementes 1 ist, so ändert sich $\Theta$ nicht, wenn man $\alpha_1, \alpha_2 \ldots \alpha_n$ in (1) durch $\alpha_1 + h_1, \alpha_2 + h_2 \ldots \alpha_n + h_n$ ersetzt. Es wird also jedes Element $\Theta$ durch (1) mindestens so oft dargestellt, als das Element 1 durch (2).

Geben andererseits die beiden Exponentenreihen $\alpha_1, \alpha_2 \ldots \alpha_n$ und $\alpha_1', \alpha_2' \ldots \alpha_n'$ zwei Darstellungen des Elementes $\Theta$, so hat man
%CONT>>
